\documentclass[a4paper,20pt]{article}
\usepackage{ctex}
\usepackage[colorlinks,linkcolor=purple]{hyperref}
\usepackage{indentfirst}
\usepackage{graphicx} %插入图片的宏包
\usepackage{float} %设置图片浮动位置的宏包
\usepackage{subfigure} %插入多图时用子图显示的宏包
\usepackage{amssymb}
\setlength{\parindent}{2em}
\title{未央-材22寒假学习手册\\第一弹：LaTeX的学习与使用}
\author{Chloral}
\begin{document}
\maketitle
\section{前言}
以下的软件无论是配置还是应用都存在一定难度,遇到意外的错误是正常情况，不出意外才是真正的意外。
（事实上，笔者在下列每个软件的安装过程中都遇到了问题QwQ）当遇到Bug时建议去知乎、CSDN论坛、
菜鸟教程上搜索解决方法，并询问班级同学。
此外，本文涉及到的所有源代码都会存到
\href{https://github.com/ChloralTHU/ChloralTHU.github.io}
{我的GitHub(https://github.com/ChloralTHU/ChloralTHU.github.io)}
中供下载查看。祝顺利成为debugger！\par
$$ {Der \quad einfache \quad Weg \quad ist \quad immer \quad verkehrt.}\tt $$


\section{LaTeX}
$ L^AT_E\mathbb{X} $可以提供专业级的排版设计，使你的文档看起来如同印刷好的一样。(事实上，本文就是笔者使用一下午速成的LaTeX语法写成的，本文档的源代码见\href{https://github.com/ChloralTHU/ChloralTHU.github.io}{这里})。\par
\subsection{LaTeX的安装}
\href{https://zhuanlan.zhihu.com/p/166523064}{LaTeX的VScode配置}
，这个链接非常详细的介绍了LaTeX及vscode扩展的安装。
简而言之，要先安装texlive,再安装vscode的LaTeX workshop扩展。
（但是这种方法对计算机水平要求相对较高，可能遇到很多bug，下载所花的时间很长。
优点：一旦配置成功接下来的所有编辑操作都很方便。）\par 
\href{https://cn.overleaf.com/}{Overleaf在线LaTeX编辑器(个人网络请点此处)}
\href{https://overleaf.tsinghua.edu.cn//}{(校园网请点击此处)}
（新手区）Overleaf是一款在线的编辑器。它可以进行文本的编辑和代码的编译，
无需配置本地环境，对新手十分友好。
Overleaf有丰富的模板库，支持多人合作，可以定位编译文档某位置的对应源码，
并且用户可以免费查看24h内的历史版本和修改记录。\par
\href{https://mp.weixin.qq.com/s/LoDPXa9BpFCEu2z-_sKZmg}{参考：未央学习资源共享计划}
,这个链接是未央书院提供的教程，但我感觉不是特别清楚。\par
\subsection{LaTeX的使用}

\subsubsection{LaTeX的基本语法}
\%是LaTeX里的注释符号，\href{https://mp.weixin.qq.com/s/LoDPXa9BpFCEu2z-_sKZmg}{此文档的最后部分}
介绍了LaTeX的基本框架、插入图片与表格的方法和插入数学公式的方法。\par
\href{https://renli1024.github.io/2018/01/19/LaTeX/LaTeX%E6%95%B0%E5%AD%A6%E5%85%AC%E5%BC%8F/}
{（非常推荐）此文}详尽而有条理地介绍了LaTeX数学公式格式。\par
\href{https://zhuanlan.zhihu.com/p/451875135}{LaTeX数学公式}
本文介绍了LaTeX的各种括号、根式分式对数式范数绝对值的输入方式。\par
\href{https://zhuanlan.zhihu.com/p/92853481}{LaTeX的高阶使用}
本文介绍了LaTeX类似编程语言的函数定义与调用。\par
在阅读上方第一篇文章插入图片的基本格式的基础上，
\href{https://blog.csdn.net/panbaoran913/article/details/122849294}{本文}
对LaTeX的图片插入做了更详细的介绍。（缺点是没有介绍插入图片需要引入的库。）\par

\subsubsection{LaTeX制作PPT}
\href{https://www.overleaf.com/learn/latex/Beamer}{本文（英文详细版）}
以及\href{https://blog.csdn.net/CoolBoySilverBullet/article/details/120803905}{本文（中文简略版）}
是LaTeX制作PPT比较清晰易懂的入门教程。\par

\begin{figure}[H] %H为当前位置，htbp为浮动图形
    \centering %图片居中
    \includegraphics[width=0.7\textwidth]{ppttry.png} %插入图片，[]中设置图片大小，{}中是图片文件名
    \caption{LaTeX制作PPT实例} %最终文档中希望显示的图片标题，可去
    \label{Fig.main1} %用于文内引用的标签，可去
    \end{figure}
\href{https://zhuanlan.zhihu.com/p/511216478}{本文}非常详细的列举了LaTeX制作PPT的优势与各种函数，但对新手不太友好。\par
\end{document}

